\begin{summarybox}
	\textbf{\textcolor{CSummary}{一句话核心}}

	辅助角公式能把同角的正弦与余弦线性组合合并为一个正弦（或余弦）函数，从而简化研究函数的性质。

	\medskip
	\textbf{\textcolor{CSummary}{使用条件}}
	\begin{itemize}[leftmargin=*, itemsep=0.5em, topsep=0.4em]
		\item \textbf{条件 1（系数不全零）：} $a,b$ 为实数，且不同时为零。
		\item \textbf{条件 2（角度任意）：} 变量 $x$ 可为任意实数。
	\end{itemize}

	\medskip
	\textbf{\textcolor{CSummary}{关键提醒}}
	\begin{itemize}[leftmargin=*, itemsep=0.5em, topsep=0.4em]
		\item \textbf{易错点（辅助角象限）：} 确定 $\varphi$ 时不能只用 $\tan\varphi$，必须结合 $\cos\varphi$ 和 $\sin\varphi$ 的符号。
		\item \textbf{检查项（系数符号）：} 使用前先标出 $a,b$ 的正负，再代入公式，避免符号错误。
	\end{itemize}
\end{summarybox}
